% Complete discrete-time martingale arguments and source exercise constructions.

\begin{annexitem}{annex:p15.martingale-examples}{thm:p15-martingale-examples}{Every conditional identity and normalization}
\begin{proof}
For a finite sum of integrable increments, \(S_n\) is integrable
and measurable from the first \(n\) increments.
The next increment is independent of their generated sigma-algebra,
so the proved conditional independence and linearity formulas give
\[
                     E(S_{n+1}\mid\mathcal F_n)=S_n+\mu.
\]
Subtracting \((n+1)\mu\) proves the centered martingale assertion.
The two signs of \(\mu\) give the other two classes.

For the quadratic construction put
\(\eta_i=\xi_i-\mu\) and \(R_n=S_n-n\mu=c+\sum_{i=1}^n\eta_i\).
The \(\eta_i\) remain independent, have mean zero and second moment
\(\sigma^2\). Finite variance addition gives
\(ER_n^2=c^2+n\sigma^2<\infty\).
In particular \(R_n\eta_{n+1}\) is integrable by Cauchy--Schwarz.
Conditional independence and the measurable-multiplier rule now yield
\[
 \begin{aligned}
 E(R_{n+1}^2-(n+1)\sigma^2\mid\mathcal F_n)
   &=R_n^2+2R_nE(\eta_{n+1}\mid\mathcal F_n)
        +E(\eta_{n+1}^2\mid\mathcal F_n)-(n+1)\sigma^2\\
   &=R_n^2+0+\sigma^2-(n+1)\sigma^2
     =R_n^2-n\sigma^2.
 \end{aligned}
\]
All terms in this identity are integrable, as required by the definition.

For independent nonnegative mean-one \(Y_i\), finite expectation
factorization gives \(E\prod_{i=1}^nY_i=1\).
Thus every \(M_n\) is adapted and integrable, and the product
\(M_nY_{n+1}=M_{n+1}\) is integrable.
The conditional multiplier and independence rules give
\[
                       E(M_{n+1}\mid\mathcal H_n)
                          =M_nE(Y_{n+1}\mid\mathcal H_n)=M_n.
\]
This uses mutual independence of the finite family, not identical laws.

For the exponential construction, \(e^{\theta\xi_i}>0\) everywhere.
Its expectation, assumed finite, cannot be zero by the zero-integral
criterion. Thus \(Y_i=e^{\theta\xi_i}/\varphi(\theta)\) is well defined,
independent across \(i\), nonnegative and of mean one.
Its product is exactly
\[
                  \prod_{i=1}^nY_i
                       =e^{\theta(S_n-c)}/\varphi(\theta)^n.
\]
The product theorem gives initial value one and the martingale
identity for the \(Y\)-filtration.
Here \(Y_{n+1}\) is also independent of the larger increment-past
\(\mathcal F_n=\sigma(\xi_1,\ldots,\xi_n)\), and the product
is \(\mathcal F_n\)-measurable. The same conditional product
calculation therefore proves the identity for \((\mathcal F_n)\).
Multiplication by the finite constant \(e^{\theta c}\) preserves
the conditional identity and gives the other displayed process.
This verifies its initial value and every normalization factor.
\end{proof}
\end{annexitem}

\begin{annexitem}{annex:p15.martingale-conditioning}{thm:p15-martingale-conditioning}{Iteration and reduction of the filtration}
\begin{proof}
For a supermartingale the one-step assertion is its definition.
Suppose \(n>m+1\). The tower property and conditional order give
\[
 \begin{aligned}
 E(X_n\mid\mathcal F_m)
   &=E(E(X_n\mid\mathcal F_{n-1})\mid\mathcal F_m)\\
   &\leq E(X_{n-1}\mid\mathcal F_m).
 \end{aligned}
\]
Induction decreases the final index to \(m+1\), then to \(X_m\).
The case \(n=m\) follows from adaptedness.
Replacing \(X\) by \(-X\) proves the submartingale assertion,
and applying both inequalities proves the martingale equality.
Taking expectations gives the corresponding mean comparisons.

For a smaller filtration \(\mathcal G_n\subseteq\mathcal F_n\)
to which \(X\) is adapted, the tower property gives
\[
 E(X_{n+1}\mid\mathcal G_n)
       =E(E(X_{n+1}\mid\mathcal F_n)\mid\mathcal G_n).
\]
For a supermartingale, conditional order bounds this by
\(E(X_n\mid\mathcal G_n)=X_n\).
The sign-reversed and equality cases are identical with their
respective signs. Integrability is unchanged, and adaptedness
was explicitly assumed.
For the natural filtration, each \(X_i\), \(i\leq n\), is
\(\mathcal F_n\)-measurable, so
\(\sigma(X_0,\ldots,X_n)\subseteq\mathcal F_n\).
This also ensures that a nonconstant initial value is retained.
\end{proof}
\end{annexitem}

\begin{annexitem}{annex:p15.martingale-convex}{thm:p15-martingale-convex}{Conditional Jensen, truncation and maxima}
\begin{proof}
Finite convex functions on the real line are Borel measurable
by the earlier supporting-line argument. Thus \(\phi(X_n)\)
is adapted, and its integrability is a hypothesis.
For a martingale, conditional Jensen gives
\[
           E(\phi(X_{n+1})\mid\mathcal F_n)
               \geq\phi(E(X_{n+1}\mid\mathcal F_n))=\phi(X_n).
\]
If \(X\) is a submartingale and \(\phi\) is also nondecreasing,
the same calculation gives a further lower bound \(\phi(X_n)\).
The real function \(|x|^p\) is convex for \(p\geq1\), as proved
in the \(L^p\) norm argument, giving the moment special case.

The function \(x\mapsto(x-a)^+\) is convex and nondecreasing,
and is bounded in absolute value by \(|x|+|a|\).
This proves its submartingale assertion without an added moment.
If \(X\) is a supermartingale, then \(-X\) is a submartingale,
so \((-X_n+a)^+\) is a submartingale.
Its negative plus the constant \(a\) is the supermartingale
\(a-(-X_n+a)^+=X_n\wedge a\).

If \(X,Y\) are submartingales for the same filtration, their
maximum is adapted and integrable, since its absolute value is
at most \(|X_n|+|Y_n|\).
Conditional order gives
\[
 E(X_{n+1}\vee Y_{n+1}\mid\mathcal F_n)
    \geq E(X_{n+1}\mid\mathcal F_n)\vee
                 E(Y_{n+1}\mid\mathcal F_n)
    \geq X_n\vee Y_n.
\]
Taking negatives proves the corresponding minimum assertion
for supermartingales.
\end{proof}
\end{annexitem}

\begin{annexitem}{annex:p15.predictable-transform}{thm:p15-predictable-transform}{Finite integrability and each one-step inequality}
\begin{proof}
For each \(i\), let \(|H_i|\leq K_i<\infty\).
If the bound is only almost everywhere, clipping \(H_i\)
to that interval preserves its \(\mathcal F_{i-1}\)-measurability
and changes every transform only almost everywhere.
The product \(H_i\Delta X_i\) is \(\mathcal F_i\)-measurable
and satisfies
\[
                E|H_i\Delta X_i|
                    \leq K_i(E|X_i|+E|X_{i-1}|)<\infty.
\]
Thus every finite transform sum is adapted and integrable.
Write \(G_n=(H\cdot X)_n\).
Conditional linearity and the proved multiplier rule yield
\[
 E(G_{n+1}\mid\mathcal F_n)
      =G_n+H_{n+1}\bigl(E(X_{n+1}\mid\mathcal F_n)-X_n\bigr).
\]
For a martingale the bracket is zero, allowing either sign of \(H\).
For a supermartingale the bracket is nonpositive, and multiplication
by \(H_{n+1}\geq0\) preserves that sign.
The submartingale case uses the opposite sign.
No step in this proof requires \(\sup_iK_i<\infty\).
\end{proof}
\end{annexitem}

\begin{annexitem}{annex:p15.stopping-calculus}{thm:p15-stopping-calculus}{The stopping sigma-algebra and the pathwise stopped sum}
\begin{proof}
The event \(\{T\leq n\}\) is the union of \(\{T=j\}\) for \(0\leq j\leq n\).
Conversely \(\{T=n\}=\{T\leq n\}\setminus\{T\leq n-1\}\),
with the second event empty at \(n=0\).
Filtration inclusion therefore proves equivalence of the two stopping
conditions. It also makes \(\{T=\infty\}\) measurable in the ambient space.

The whole space belongs to \(\mathcal F_T\).
For a member \(A\),
\[
 A^c\cap\{T\leq n\}
     =\{T\leq n\}\setminus(A\cap\{T\leq n\})\in\mathcal F_n.
\]
Countable unions commute with the same intersection, proving that
\(\mathcal F_T\) is a sigma-algebra.
If \(S\leq T\), then for \(A\in\mathcal F_S\),
\[
 A\cap\{T\leq n\}
       =(A\cap\{S\leq n\})\cap\{T\leq n\}\in\mathcal F_n,
\]
so \(\mathcal F_S\subseteq\mathcal F_T\).

For \(T\leq N<\infty\), the variable
\(X_T=\sum_{j=0}^NX_j\mathbf1_{\{T=j\}}\) is integrable:
its absolute expectation is at most \(\sum_{j=0}^NE|X_j|\).
For any Borel set \(B\),
\[
 \{X_T\in B\}\cap\{T\leq n\}
      =\bigcup_{j=0}^{\min(n,N)}
                     (\{T=j\}\cap\{X_j\in B\})\in\mathcal F_n.
\]
Hence \(X_T\) is \(\mathcal F_T\)-measurable.
If a deterministic bound holds only almost everywhere, use
\(X_{T\wedge N}\) as the version on the exceptional set;
it is measurable for \(\mathcal F_{T\wedge N}\subseteq\mathcal F_T\).

For an arbitrary \(T\), the coefficient \(\mathbf1_{\{T\geq i\}}\)
is predictable, because
\(\{T\geq i\}=\{T\leq i-1\}^c\in\mathcal F_{i-1}\).
On each outcome, the finite sum
\(\sum_{i=1}^n\mathbf1_{\{T\geq i\}}\Delta X_i\)
telescopes exactly to \(X_{T\wedge n}-X_0\).
The predictable-transform theorem supplies its martingale class.
The process constantly equal in time to the integrable
\(\mathcal F_0\)-measurable variable \(X_0\) is a martingale.
Adding it preserves the relevant conditional equality or inequality.
This proves the stopped-process result, including \(T=\infty\).
\end{proof}
\end{annexitem}

\begin{annexitem}{annex:p15.bounded-optional-sampling}{thm:p15-bounded-optional-sampling}{The exact stopped-information event calculation}
\begin{proof}
The stopped values \(X_S,X_T\) are integrable, and \(X_S\) is
\(\mathcal F_S\)-measurable by the stopping-calculus theorem.
For \(A\in\mathcal F_S\), the pathwise difference gives
\[
       \mathbf1_A(X_T-X_S)
              =\sum_{i=1}^N
                     \mathbf1_{A\cap\{S<i\}\cap\{T\geq i\}}\Delta X_i.
\]
Each coefficient is \(\mathcal F_{i-1}\)-measurable:
\(A\cap\{S<i\}=A\cap\{S\leq i-1\}\) has that measurability
by the definition of \(\mathcal F_S\), and \(\{T\geq i\}\)
has it by predictability.
Each term is integrable since the coefficient is bounded.
For a supermartingale, its expectation is
\[
 E\left[
 \mathbf1_{A\cap\{S<i\}\cap\{T\geq i\}}
       \bigl(E(X_i\mid\mathcal F_{i-1})-X_{i-1}\bigr)\right]\leq0.
\]
Summing yields \(\int_AX_T\,dP\leq\int_AX_S\,dP\) for every
\(A\in\mathcal F_S\).
The bounded-level conditional comparison argument then proves
\(E(X_T\mid\mathcal F_S)\leq X_S\).
For a martingale every term is zero, giving the defining identity
and equality by uniqueness.
The submartingale case reverses the inequalities.
If the bounds hold almost everywhere, the same calculation is
unchanged after assigning the stopped versions \(X_{S\wedge N}\)
and \(X_{T\wedge N}\) on the exceptional set.
Finally take \(S=0\) and then expectations for the bounded-time
martingale mean identity.
\end{proof}
\end{annexitem}

\begin{annexitem}{annex:p15.upcrossing-inequality}{thm:p15-upcrossing-inequality}{Alternating hitting times, greedy counting and clipping}
\begin{proof}
For the first purchase,
\(\{T_1\leq n\}=\bigcup_{m=0}^n\{X_m\leq a\}\in\mathcal F_n\).
If an earlier alternating time \(T_j\) is a stopping time, then
the event that its next threshold is reached by \(n\) is the
finite union over \(0\leq k<m\leq n\) of
\(\{T_j=k\}\cap\{X_m\text{ satisfies the next threshold}\}\).
Each member is \(\mathcal F_n\)-measurable.
This proves recursively that all \(T_j\), \(j\geq1\), are stopping times.
The rule for an infinite preceding time is exactly the empty-search rule.

The holding event for interval \(j\) at time \(m\geq1\) is
\[
       \{T_{2j-1}\leq m-1\}\cap\{T_{2j}>m-1\}\in\mathcal F_{m-1}.
\]
Their countable union is measurable there.
At most one such interval is active, because the successive finite
times strictly increase and \(a<b\).
Thus \(H_m\) is predictable and takes only the values zero and one.

For a deterministic finite path, each completed greedy pair
\((T_{2j-1},T_{2j})\) is an admissible low-to-high crossing.
Conversely take any \(k\) chronologically disjoint pairs
\(s_1<t_1<s_2<t_2<\cdots<s_k<t_k\leq n\)
with \(X_{s_j}\leq a,X_{t_j}\geq b\).
The first greedy purchase is no later than \(s_1\), and its sale
is no later than \(t_1\), since those times are available candidates.
Inductively the next proposed purchase \(s_j\) and sale \(t_j\)
are still later than the preceding greedy times.
Hence \(T_{2j-1}\leq s_j,T_{2j}\leq t_j\) for all \(j\).
The greedy count is at least \(k\), proving maximality.
There can be only finitely many pairs at finite horizon.

Now let \(X\) be a submartingale and put
\(Y_m=a+(X_m-a)^+=X_m\vee a\).
It is an integrable submartingale by the convex-transform theorem.
The conditions \(Y_m\leq a\) and \(Y_m\geq b\) are equivalent
to \(X_m\leq a\) and \(X_m\geq b\), respectively.
Thus it has exactly the same alternating times and count.
On each completed holding interval the transform gain is
\[
                    Y_{T_{2j}}-Y_{T_{2j-1}}\geq b-a,
\]
because the purchase value is exactly \(a\).
If there is an uncompleted holding interval at time \(n\),
its contribution is \(Y_n-a\geq0\), including zero for a purchase
made at \(n\). Summing the increments over the disjoint holding
intervals gives the pathwise bound
\[
                              (b-a)U_n(a,b)\leq(H\cdot Y)_n.
\]

Let \(K_m=1-H_m\), another bounded nonnegative predictable sequence.
The finite telescoping identity is
\[
                         Y_n-Y_0=(H\cdot Y)_n+(K\cdot Y)_n.
\]
The latter transform is a submartingale starting at zero, so
\(E(K\cdot Y)_n\geq0\).
Taking expectations in the two displayed inequalities gives
\[
 (b-a)EU_n(a,b)\leq E(Y_n-Y_0)
                  =E[(X_n-a)^+]-E[(X_0-a)^+].
\]
Every finite-horizon transform is integrable by its proved bound.
This establishes the source inequality with its actual clipped process.
\end{proof}
\end{annexitem}

\begin{annexitem}{annex:p15.martingale-convergence}{thm:p15-martingale-convergence}{All rational crossings and both infinite endpoints}
\begin{proof}
Assume \(X\) is a submartingale and \(C=\sup_nEX_n^+<\infty\).
Since \((x-a)^+\leq x^++|a|\), the upcrossing inequality gives
\[
                         EU_n(a,b)\leq\frac{C+|a|}{b-a}.
\]
The counts increase to \(U_\infty(a,b)\).
Monotone convergence gives the same finite upper bound for its
expectation, hence \(U_\infty(a,b)<\infty\) almost surely.
Apply this to the countable family of rational \(a<b\) and
remove their union of exceptional null sets.

If on one remaining path
\(\liminf_nX_n<\limsup_nX_n\), there are rationals
\(a<b\) strictly between those extended limits.
There are arbitrarily late values below \(a\), and, after every
such time, later values above \(b\).
Selecting them alternately constructs infinitely many chronological
crossings, contradicting the finite count.
Thus the extended-real limit exists almost surely.
The limit is measurable, since limsup and liminf of a measurable
sequence are measurable.

Let that extended limit be \(L\). Fatou's lemma gives
\[
                          EL^+\leq\liminf_nEX_n^+\leq C.
\]
So \(P(L=+\infty)=0\).
The submartingale expectation comparison also gives
\[
                 EX_n^-=EX_n^+-EX_n\leq C-EX_0 .
\]
Another application of Fatou yields
\(EL^-\leq C-EX_0<\infty\), excluding \(L=-\infty\).
Assign zero on the measurable null set where a finite limit
does not exist, obtaining a finite measurable \(X_\infty\).
The two bounds prove integrability.
Fatou applied to \(|X_n|\) gives the stated sharper absolute
moment bound.

For a supermartingale with bounded negative-part means, apply
the result to \(-X\).
For a nonnegative supermartingale this negative part is zero.
The limit is nonnegative and integrable, and Fatou combined
with \(EX_n\leq EX_0\) gives \(EX_\infty\leq EX_0\).
For each fixed \(n\), conditional Fatou for the nonnegative
sequence \(X_{n+k}\) gives the further relation
\[
 E(X_\infty\mid\mathcal F_n)
       \leq\liminf_kE(X_{n+k}\mid\mathcal F_n)\leq X_n.
\]
An \(L^1\)-bounded martingale in particular has bounded positive
means and so satisfies the theorem.
None of these steps gives convergence of the first absolute
moments; the following explicit source example shows why that
additional conclusion cannot be added.
\end{proof}
\end{annexitem}

\begin{annexitem}{annex:p15.martingale-source-counterexamples}{thm:p15-martingale-source-counterexamples}{Absorption and an independently driven intermittent process}
\begin{proof}
\paragraph{The stopped symmetric walk.}
The independent-sign construction gives a martingale \(S\)
starting at 1. Its first hitting time \(T\) of zero is a stopping
time, since occurrence by \(n\) is a finite union of observed events.
Until zero is reached, a walk with integer values and increments
\(\pm1\) cannot become negative: a first passage from a positive
integer to a nonpositive integer must pass through zero.
Thus the stopped process \(X_n=S_{T\wedge n}\) is nonnegative.
It is a martingale by the stopping theorem, so \(EX_n=1\) at
every finite time.
The nonnegative-supermartingale convergence theorem gives a
finite almost-sure limit.

A convergent sequence of integers is eventually constant:
once its terms lie in an interval of length less than one,
they must all equal its single possible integer.
On \(\{T=\infty\}\), however, every successive increment
of \(X\) is still \(+1\) or \(-1\), so it is never eventually
constant. This event must therefore be null.
Hence \(T<\infty\) almost surely and the limit is zero.
Nevertheless \(E|X_n-0|=EX_n=1\), so there is no \(L^1\)
convergence. Recurrence of the walk was derived here, not
assumed as a separate theorem.

\paragraph{The intermittent martingale.}
The countable-product construction supplies the independent
Bernoulli and sign variables in the statement.
Take the filtration generated by their first \(k\) pairs.
The recursive \(X_k\) is adapted, and
\(|X_k|\leq1+k|X_{k-1}|\) proves finite integrability by induction.
Write its recursion as
\[
 X_k=\mathbf1_{\{X_{k-1}=0\}}B_k\varepsilon_k
           +\mathbf1_{\{X_{k-1}\ne0\}}kX_{k-1}B_k.
\]
The fresh variables are independent of the preceding filtration.
Their means are \(EB_k=1/k\) and \(E\varepsilon_k=0\).
Conditional factorization and the integrable multiplier rule
therefore give \(E(X_k\mid\mathcal F_{k-1})=X_{k-1}\).

Every time \(B_k=1\), the new value is a nonzero integer, and
every time \(B_k=0\), the new value is zero.
Thus \(\{X_k\ne0\}=\{B_k=1\}\) exactly, giving probability
\(1/k\) and convergence to zero in probability.
The events \(B_k=1\) are independent and their probabilities
have divergent sum. The second Borel--Cantelli lemma gives
nonzero values infinitely often almost surely.
The events \(B_k=0\), \(k\geq2\), are also independent and
have divergent sum, so zero values occur infinitely often too.
All nonzero values have absolute value at least one.
These two facts exclude every finite or infinite path limit.

For clarity, its absolute moments can be computed rather than
presuming they are uniformly bounded:
\[
                 E|X_k|=E|X_{k-1}|+\frac{P(X_{k-1}=0)}k.
\]
Here \(E|X_1|=1\) and, for \(k\geq2\),
\(P(X_{k-1}=0)=1-1/(k-1)\).
Summing gives
\(E|X_k|=\sum_{j=1}^k1/j-1+1/k\), which diverges.
Thus the construction does not meet the convergence theorem's
uniform moment condition.
\end{proof}
\end{annexitem}

\begin{annexitem}{annex:p15.localized-martingale-convergence}{thm:p15-localized-martingale-convergence}{A stopping level and its integrable overshoot}
\begin{proof}
Let \(D=\sup_{n\geq1}(\Delta X_n)^+\), which is measurable,
nonnegative and integrable by hypothesis.
For each positive integer \(K\), let
\(T_K=\inf\{n\geq0:X_n>K\}\), a stopping time.
Before that time \(X_n\leq K\).
If \(1\leq T_K<\infty\), then
\[
                    X_{T_K}\leq X_{T_K-1}+D\leq K+D.
\]
If \(T_K=0\), the stopped value is \(X_0\).
Thus for every \(n\),
\[
                    (X_{T_K\wedge n})^+\leq X_0^++K+D
                                           \quad\text{almost everywhere}.
\]
The stopped process is a submartingale and has uniformly bounded
positive-part expectations. The proved convergence theorem makes
it converge almost surely to a finite limit.
Remove the union of these null exceptions over all positive integer \(K\).

On any remaining path with \(\sup_nX_n<\infty\), choose an integer
\(K\) exceeding that supremum. Then \(T_K=\infty\) on the path,
and the convergent stopped sequence is exactly the original sequence.
The hypothesis makes these paths have full probability.
Hence \(X_n\) has a finite almost-sure limit, as claimed.
\end{proof}
\end{annexitem}

\begin{annexitem}{annex:p15.nontrivial-product-limit}{thm:p15-nontrivial-product-limit}{The product limit and all logarithmic cases}
\begin{proof}
The product is a nonnegative mean-one martingale by the earlier
construction, so it converges almost surely to a finite \(M_\infty\).
On \(\{M_\infty>0\}\), no product factor can have been zero,
and the convergent positive sequence satisfies
\[
                          Y_n=M_n/M_{n-1}\longrightarrow1.
\]
Since \(P(Y_1\ne1)>0\), some \(\varepsilon>0\) has
\(P(|Y_1-1|>\varepsilon)>0\).
The corresponding events for the i.i.d. factors are independent
with a divergent sum of probabilities.
Borel--Cantelli makes them occur infinitely often almost surely.
Thus \(P(Y_n\to1)=0\), and \(P(M_\infty>0)=0\).
The nonnegative limit must be zero.

For the logarithmic rate, define \(\log0=-\infty\).
For \(y>0\), \((\log y)^+\leq y\), so its expectation is finite.
Therefore \(c=E\log Y_1\) is well defined in \([-\infty,\infty)\).
If \(P(Y_1=0)=p_0>0\), independence gives probability
\((1-p_0)^n\) for avoiding a zero through time \(n\).
A zero occurs at a finite time almost surely, and the product
and its logarithm are thereafter zero and \(-\infty\), respectively.
In this case \(c=-\infty\) and the asserted rate follows.

Suppose now the factors are positive almost surely.
After removing the countable union of their null zero events,
all logarithms are finite real numbers and
\(\log M_n=\sum_{i=1}^n\log Y_i\).
If the negative logarithmic part also has finite expectation,
the integrable i.i.d. strong law in \cref{thm:p15-iid-series-strong-law} gives
\(n^{-1}\log M_n\to c\).
If its expectation is infinite, let
\[
                         Z_i^{(K)}=\max(\log Y_i,-K)
                                           \qquad(K=1,2,\ldots).
\]
Each sequence is i.i.d. and integrable.
Its mean \(c_K=E(\log Y_1)^+-E\min((\log Y_1)^-,K)\)
decreases to \(-\infty\) by monotone convergence.
For every \(K\), the strong law and the pointwise inequality
\(\log Y_i\leq Z_i^{(K)}\) give
\[
                    \limsup_n\frac1n\log M_n\leq c_K
                                         \quad\text{almost surely}.
\]
Remove their countable exceptional union and let \(K\to\infty\).
This proves convergence to \(-\infty\).

Finally, in the finite-logarithm case, the elementary inequality
\(\log y\leq y-1\), strict except at \(y=1\), implies
\[
                       E(Y_1-1-\log Y_1)=-c>0.
\]
The strict positivity follows from the zero-integral criterion
and \(P(Y_1\ne1)>0\); all terms are integrable in this case.
Thus \(c<0\). In the remaining cases \(c=-\infty<0\) already.
This proves the entire extended range and rate statement.
\end{proof}
\end{annexitem}

\begin{annexitem}{annex:p15.summable-product-compensation}{thm:p15-summable-product-compensation}{A positive product limit and the normalized supermartingale}
\begin{proof}
First let \(y_n>-1\) and \(\sum_n|y_n|<\infty\).
Then \(y_n\to0\), so eventually \(|y_n|\leq1/2\).
On this interval, integrating the derivative \(1/(1+t)\)
between zero and \(y_n\) gives
\[
                           |\log(1+y_n)|\leq2|y_n|.
\]
The tail logarithm series is absolutely convergent.
Its finitely many initial terms are finite because every factor
is strictly positive. Hence the entire logarithm series converges
to a finite real number, and exponentiating gives a finite
strictly positive product limit.

In the stochastic assertion, \(A_n\) is finite, positive and
\(\mathcal F_{n-1}\)-measurable for \(n\geq1\), and \(A_n\geq1\).
The variable \(Z_n=X_n/A_n\) is adapted and integrable since
\(0\leq Z_n\leq X_n\).
The reciprocal \(1/A_{n+1}\) is \(\mathcal F_n\)-measurable
and bounded by one. Thus the conditional multiplier rule gives
\[
 \begin{aligned}
 E(Z_{n+1}\mid\mathcal F_n)
   &=\frac{E(X_{n+1}\mid\mathcal F_n)}{A_n(1+Y_n)}\\
   &\leq\frac{(1+Y_n)X_n}{A_n(1+Y_n)}=Z_n.
 \end{aligned}
\]
The product \((1+Y_n)X_n\) need not have a finite expectation:
the argument only multiplies its given pointwise upper bound
by a finite reciprocal, while the conditional identity is
applied to the integrable \(X_{n+1}\).
Hence \(Z\) is a nonnegative supermartingale and has a finite
almost-sure limit.
On the full-probability event \(\sum_nY_n<\infty\), the
deterministic product result, with its index shifted by one,
gives \(A_n\to A_\infty\in[1,\infty)\), finite.
Therefore \(X_n=A_nZ_n\to A_\infty Z_\infty\), a finite limit.
No expectation bound for this unnormalized product is used.
\end{proof}
\end{annexitem}

\begin{annexitem}{annex:p15.supermartingale-switching}{thm:p15-supermartingale-switching}{Both discrete switching conventions}
\begin{proof}
Both processes are adapted, since every displayed stopping
event belongs to \(\mathcal F_n\), and integrable because
their absolute values are at most \(|X^1_n|+|X^2_n|\).
The comparison at the switch gives \(Y_n\leq Z_n\) almost
everywhere: they differ only on \(\{T=n\}\), where their
values are respectively \(X^2_n,X^1_n\).
For the next delayed value, the exact integer inequalities give
\[
        Z_{n+1}=X^1_{n+1}\mathbf1_{\{T>n\}}
                         +X^2_{n+1}\mathbf1_{\{T\leq n\}}.
\]
The two indicators are \(\mathcal F_n\)-measurable.
Conditioning and using the two supermartingale inequalities yields
\[
                    E(Z_{n+1}\mid\mathcal F_n)
                       \leq X^1_n\mathbf1_{\{T>n\}}
                              +X^2_n\mathbf1_{\{T\leq n\}}=Y_n.
\]
Since \(Y_{n+1}\leq Z_{n+1}\), conditional order gives
\(E(Y_{n+1}\mid\mathcal F_n)\leq Y_n\).
The inequality \(Y_n\leq Z_n\) gives
\(E(Z_{n+1}\mid\mathcal F_n)\leq Z_n\).
This proves both processes are supermartingales.
The assumption at \(T\) is used only on finite switch events,
so no \(X^j_\infty\) value is required.
\end{proof}
\end{annexitem}

\begin{annexitem}{annex:p15.dubins-upcrossing}{thm:p15-dubins-upcrossing}{The finite switching sequence behind Dubins' estimate}
\begin{proof}
Use the alternating times \(T_j\) for \(0<a<b\), already
proved to be stopping times. Define the nonnegative process
\(Y\) by value 1 before \(T_1\) and, for \(j\geq1\),
\[
 Y_n=
 \begin{cases}
  (b/a)^{j-1}X_n/a,&T_{2j-1}\leq n<T_{2j},\\
  (b/a)^j,&T_{2j}\leq n<T_{2j+1}.
 \end{cases}
\]
At any finite \(n\) there have been only finitely many alternating
times, so exactly one of these cases applies.
In particular
\[
                              Y_0=\min(X_0/a,1).
\]
If \(T_1=0\) this is \(X_0/a\), and otherwise \(X_0>a\)
and the value is one.

We prove the source's stopped-process assertion using finite
switching, avoiding an assumption about an infinite composition.
Start with the constant process 1.
At \(T_1\), switch to \(X_n/a\).
Its value at that time is at most the old value because
\(X_{T_1}\leq a\).
The switching principle makes the new process a supermartingale.
At \(T_2\), switch to the constant \(b/a\); the preceding
scaled process has value at least \(b/a\) there.
Continue: at \(T_{2j-1}\), switch from the constant
\((b/a)^{j-1}\) to \((b/a)^{j-1}X_n/a\), a downward
switch because \(X_{T_{2j-1}}\leq a\).
At \(T_{2j}\), switch from that scaled process to
\((b/a)^j\), again downward because \(X_{T_{2j}}\geq b\).
Every new component is an integrable supermartingale, being
a deterministic finite multiple of \(X\) or a constant.
Induction proves that the process after any finite number
of switches is a nonnegative supermartingale.

After the first \(j\) switches it agrees with \(Y\) through
time \(T_j\). Stopping it at \(T_j\) therefore proves that
\[
                              Z_n^{(j)}=Y_{n\wedge T_j}
\]
is a nonnegative supermartingale.
Its initial value is \(Y_0\) for every \(j\geq1\).
This also justifies its finite-time integrability when some
alternating times are infinite.

Fix \(k\geq1\). On \(\{T_{2k}\leq n\}\),
\(Z_n^{(2k)}=(b/a)^k\), and elsewhere it is nonnegative.
Consequently
\[
             (b/a)^kP(T_{2k}\leq n)
                    \leq EZ_n^{(2k)}\leq EY_0.
\]
The increasing union of these events is
\(\{T_{2k}<\infty\}=\{U_\infty(a,b)\geq k\}\).
Taking its probability limit gives the desired estimate.
Equivalently, the source's limiting step follows by Fatou:
on this event the stopped sequence is eventually the constant
\((b/a)^k\), while its nonnegative liminf elsewhere gives no
negative contribution. The same expectation bound follows.
This proves both requested stopped supermartingales and the bound.
\end{proof}
\end{annexitem}

\begin{annexitem}{annex:p15.martingale-boundary-calculations}{thm:p15-martingale-boundary-calculations}{The exact finite gains and the failed hypotheses}
\begin{proof}
For the fair walk, \(H_n=\xi_n\) gives
\(H_n\Delta S_n=\xi_n^2=1\), hence gain \(n\).
If \(H_n\) were \(\mathcal F_{n-1}\)-measurable, its
conditional expectation there would equal itself.
Independence instead gives \(E(\xi_n\mid\mathcal F_{n-1})=0\),
contradicting its values \(\pm1\).
Thus predictability fails, and the increasing deterministic
gain process is not a martingale.

For the reset/doubling strategy, induction gives
\(1\leq H_n\leq2^{n-1}\), and the recursion uses only
the previous increments, so it is predictable.
The first-success time is a stopping time, and
\(P(T>n)=(1-p)^n\), which tends to zero for \(p>0\).
Before the first success, the stakes are \(1,2,\ldots,2^{n-1}\).
After \(k\) consecutive losses and one success, the gain is
\[
                    -(1+2+\cdots+2^{k-1})+2^k=1.
\]
For the source's displayed four losses followed by a success,
the finite table is
\[
 \begin{array}{c|rrrrr}
 H_n&1&2&4&8&16\\
 \xi_n&-1&-1&-1&-1&1\\
 (H\cdot S)_n&-1&-3&-7&-15&1 .
 \end{array}
\]
Stopping at that first success gives exactly the two cases
for \(G_n\) in the statement. Its expectation is
\[
        [1-(1-p)^n]+(1-2^n)(1-p)^n=1-[2(1-p)]^n.
\]
Also \(|G_n-1|=2^n\mathbf1_{\{T>n\}}\).
At \(p=1/2\), its expectation is one for every \(n\),
despite \(G_n\to1\) almost surely.
For \(0<p<1\), \(P(T>n)>0\) at every finite \(n\);
there is no deterministic time bound. For \(p=1\),
the exceptional deterministic case is \(T=1\).
Thus none of the fair-case calculations contradict bounded
optional sampling.

For the rare-jump example, the countable-product construction
supplies independent coordinates with the specified two-point laws.
Their means are
\[
                 (n^2-1)/n^2-(1-1/n^2)=0,
\]
and their first absolute moments are
\(2(1-1/n^2)\leq2\), including the constant zero first variable.
The sums are adapted, integrable at every finite time, and
have conditionally zero mean increments, so are a martingale.
The exceptional-value probabilities sum to a finite number:
for example the dyadic block from \(2^j\) to \(2^{j+1}-1\)
contributes at most \(2^{-j}\).
Borel--Cantelli implies that only finitely many such jumps occur.
Thereafter every increment is \(-1\), so on each such path
the partial sum equals a finite path-dependent constant minus \(n\).
This proves both asserted limits.

Finally \(-1/(n+1)\) is a deterministic increasing sequence,
so satisfies the submartingale inequality for any filtration.
Its square is deterministic and decreasing, so is a
supermartingale. The square function is convex but is
decreasing on this negative range, explaining the failure
of the nondecreasing-transform hypothesis in this example.
\end{proof}
\end{annexitem}
